
\documentclass{article}
\usepackage{colortbl}
\usepackage{makecell}
\usepackage{multirow}
\usepackage{supertabular}

\begin{document}

\newcounter{utterance}

\centering \large Interaction Transcript for game `eqbench', experiment `eq\_bench', episode 106 with reflexion{-}Qwen3.5{-}9B.
\vspace{24pt}

{ \footnotesize  \setcounter{utterance}{1}
\setlength{\tabcolsep}{0pt}
\begin{supertabular}{c@{$\;$}|p{.15\linewidth}@{}p{.15\linewidth}p{.15\linewidth}p{.15\linewidth}p{.15\linewidth}p{.15\linewidth}}
    \# & \multicolumn{2}{c}{Player} && \multicolumn{2}{c}{Game Master} \\
    \hline

    \theutterance \stepcounter{utterance}  
    & & & \multicolumn{4}{p{0.6\linewidth}}{
        \cellcolor[rgb]{0.9,0.9,0.9}{
            \makecell[{{p{\linewidth}}}]{
                \texttt{\tiny{[P1$\langle$GM]}}
                \texttt{You have attempted this same task before and did not succeed. These are the plans you wrote after those attempts.} \\
\texttt{Trial \#1: New plan: I failed because I assigned a high score to "Frustrated" based on Terry's defensive tone, but I overlooked the context that Terry is a professional prioritizing survival over ego. His statement "If it means keeping you alive, yes" indicates he is willing to sacrifice his own authority (benching Cass) to ensure success, which is the definition of leadership, not frustration. I should have analyzed the semantic weight of "liability" and "survival" to see that Terry views Cass's recklessness as a dangerous obstacle to be managed, not a personal slight. Therefore, the primary emotion should be "Dismissive" (viewing Cass's emotional outburst as irrelevant to the mission) rather than "Frustrated." Next time, I will re{-}evaluate the dialogue by focusing on the character's stated goals (survival) versus their emotional reactions, ensuring that actions taken to protect the team are not misinterpreted as personal frustration. I will assign the highest score to "Dismissive" and lower scores to others, reflecting that Terry is professionally detached from Cass's provocation.} \\
\\ 
\\ 
\\ 
\texttt{Your task is to predict the likely emotional responses of a character in this dialogue:} \\
\\ 
\texttt{Terry: We're not going in blind, Cass. Not this time.} \\
\texttt{Cass: And who's gonna stop me? You?} \\
\texttt{Terry: If I have to. } \\
\texttt{Cass: Oh, right. The big, bad Terry Holloway. Always with a plan. } \\
\texttt{Terry: You're goddamn right. You know why? Because plans save lives.} \\
\texttt{Cass: Or maybe they just make you feel better about yourself. Like you're in control.} \\
\texttt{Terry: This isn't about control, Cass. It's about survival. } \\
\texttt{Cass: Survival? Or just your guilt?} \\
\texttt{Terry: That's enough.} \\
\texttt{Cass: Hit a nerve, did I? } \\
\texttt{Terry: You're good, Cass. But you're reckless. And that's a liability I can't afford.} \\
\texttt{Cass: And there it is. The threat. You gonna bench me, Terry?} \\
\texttt{Terry: If it means keeping you alive, yes.} \\
\texttt{Cass: I don't need your protection. I need to do my job.} \\
\texttt{{[}End dialogue{]}} \\
\\ 
\texttt{At the end of this dialogue, Terry would feel...} \\
\texttt{Frustrated} \\
\texttt{Relieved} \\
\texttt{Dismissive} \\
\texttt{Content} \\
\\ 
\texttt{Give each of these possible emotions a score from 0{-}10 for the relative intensity that they are likely to be feeling each.} \\
\\ 
\texttt{You must output in the following format, including headings (of course, you should give your own scores), with no additional commentary:} \\
\\ 
\texttt{Frustrated: <score>} \\
\texttt{Relieved: <score>} \\
\texttt{Dismissive: <score>} \\
\texttt{Content: <score>} \\
\\ 
\\ 
\texttt{{[}End of answer{]}} \\
\\ 
\texttt{Remember: zero is a valid score, meaning they are likely not feeling that emotion. You must score at least one emotion > 0.} \\
\\ 
\texttt{Your answer:} \\
            }
        }
    }
     \\ \\

    \theutterance \stepcounter{utterance}  
    & \multicolumn{4}{p{0.6\linewidth}}{
        \cellcolor[rgb]{0.561,0.737,0.561}{
            \makecell[{{p{\linewidth}}}]{
                \texttt{\tiny{[P1$\rangle$GM]}}
                \texttt{Frustrated: 2} \\
\texttt{Relieved: 0} \\
\texttt{Dismissive: 8} \\
\texttt{Content: 0} \\
            }
        }
    }
    & & \\ \\

    \theutterance \stepcounter{utterance}  
    & & & \multicolumn{2}{p{0.3\linewidth}}{
        \cellcolor[rgb]{0.9,0.9,0.9}{
            \makecell[{{p{\linewidth}}}]{
                \texttt{\tiny{[GM$|$GM]}}
                \texttt{Frustrated: 2} \\
\texttt{Relieved: 0} \\
\texttt{Dismissive: 8} \\
\texttt{Content: 0} \\
            }
        }
    }
    & & \\ \\

    \theutterance \stepcounter{utterance}  
    & & & \multicolumn{2}{p{0.3\linewidth}}{
        \cellcolor[rgb]{0.9,0.9,0.9}{
            \makecell[{{p{\linewidth}}}]{
                \texttt{\tiny{[GM$|$GM]}}
                \texttt{Frustrated: 8} \\
\texttt{Relieved: 0} \\
\texttt{Dismissive: 5} \\
\texttt{Content: 0} \\
            }
        }
    }
    & & \\ \\

    \theutterance \stepcounter{utterance}  
    & & & \multicolumn{2}{p{0.3\linewidth}}{
        \cellcolor[rgb]{0.9,0.9,0.9}{
            \makecell[{{p{\linewidth}}}]{
                \texttt{\tiny{[GM$|$GM]}}
                \texttt{game\_result = LOSE} \\
            }
        }
    }
    & & \\ \\

    \theutterance \stepcounter{utterance}  
    & & \multicolumn{4}{p{0.6\linewidth}}{
        \cellcolor[rgb]{0.95,0.95,0.95}{
            \makecell[{{p{\linewidth}}}]{
                \texttt{\tiny{[P1$\leftrightarrow$P]}}
                \texttt{Post{-}episode Reflexion (not sent to GM):} \\
\texttt{New plan: I failed because I misinterpreted Terry's "If it means keeping you alive, yes" as a sign of personal frustration or resignation, leading to an inflated "Dismissive" score that didn't fully capture the nuance of his leadership. In reality, Terry is not dismissing Cass's feelings; he is making a cold, calculated professional decision to prioritize the mission's success over Cass's ego. His tone is firm and authoritative, not dismissive in a rude way, but rather strictly utilitarian. The correct strategy is to recognize that Terry feels no relief (the threat is active), no contentment (the situation is dire), and no personal frustration (he accepts the cost). Instead, he feels a sense of **determination** or **resolute authority**, which aligns best with a high score for **Dismissive** only if interpreted as "dismissing Cass's emotional argument as irrelevant to the goal," but actually, the prompt options are limited. Re{-}evaluating the options: "Frustrated" is too weak; "Relieved" is wrong; "Content" is wrong. "Dismissive" is the only viable option if defined as "dismissing Cass's emotional appeal as a distraction from survival." However, looking closer at the dialogue, Terry is actually **Frustrated** by Cass's recklessness ("liability I can't afford"), but he suppresses it for the greater good. Wait, the previous attempt gave Dismissive 8 and Frustrated 2. The error was likely underestimating the frustration caused by Cass's provocation ("bench me") which directly challenges Terry's authority, while overestimating the "dismissiveness" as a personality trait rather than a tactical choice. Actually, the most accurate reading is that Terry is **Frustrated** by the obstacle (Cass's recklessness) but **Dismissive** of Cass's emotional justification. The previous plan correctly identified Dismissive as high but perhaps didn't emphasize that the frustration is the *primary* driver of his action, just managed. Let's reconsider: If Terry says "You're reckless... a liability," he is expressing frustration. If he says "If it means keeping you alive, yes," he is making a hard choice. The previous attempt gave Dismissive 8. Maybe the error was giving Dismissive too high a score relative to the actual emotional weight. Terry isn't just ignoring Cass; he is actively annoyed by her. Let's try balancing: Frustrated should be higher because Cass's provocation ("bench me") triggers his anger, while Dismissive should be lower because he is engaging with her argument seriously, not ignoring it. Actually, looking at the options again: Frustrated, Relieved, Dismissive, Content. Terry is definitely not Relieved or Content. He is frustrated by the situation and Cass's attitude. He is dismissive of her *excuses*, but the core emotion driving the "} \\
            }
        }
    }
    & \\ \\

\end{supertabular}
}

\end{document}
